% GENERATED FILE. Edit canonical lessons, then run npm run generate:books.
% canonical-source-hash: fnv1a64:e06943020af633ed
% canonical-lessons: BN-C440-bismoy, BN-C440-koutuhol, BN-C440-mugdho, BN-C440-obak-hooya, BN-C440-chomkano

\chapter{Surprise and Wonder}
\label{ch:bn-surprise-and-wonder}
\hlchaptermodality{440}

\begin{chapteropening}
\textbf{By the end of this chapter:} \emph{I can say that something surprised or amazed me.}

Complete the last lesson of chapter 440: I can say that something surprised or amazed me.
\end{chapteropening}

\section[\texorpdfstring{\bn{বিস্ময়} (bismôy)}{বিস্ময় (bismôy)}]{\bn{বিস্ময়} (bismôy) — wonder; amazement}

\label{lesson:BN-C440-bismoy}



\begin{quote}
Before the new one: say the Bengali for to enjoy, then the Bengali for to get carried away.
\end{quote}

\subsection*{You'll want to know: \bn{বিস্ময়}}
\textbf{\bn{বিস্ময়}} — \emph{bismôy} — \textquotedblleft{}wonder; amazement\textquotedblright{}.

\begin{quote}
\bn{শিশুটি} \bn{বিস্ময়ে} \bn{চাঁদটা} \bn{দেখছিল।}
\end{quote}

\begin{quote}
\emph{shishuṭi bismôye chā̃dṭā dekhchhilô.}
\end{quote}

\begin{quote}
The child was gazing at the moon in wonder.
\end{quote}

\begin{cousinweb}
\textbf{\bn{বিস্ময়}} (\emph{bismôy}) is a tatsama, Sanskrit \emph{vismaya}, \textquotedblleft{}wonder, surprise, amazement\textquotedblright{} (Monier-Williams).
\end{cousinweb}

\begin{grammarlens}[title={What you've built}]
One more word for describing what you lived and felt.
\end{grammarlens}

\subsection*{Your turn}
\begin{itemize}
\raggedright
  \item \emph{Say it:} \emph{bismôy}
  \item \emph{Say it:} \emph{bismôy}, once more
  \item \emph{Say it:} \emph{shishuṭi bismôye chā̃dṭā dekhchhilô.} — The child was gazing at the moon in wonder.
  \item \emph{Recall:} say the Bengali for to enjoy, then the Bengali for to get carried away, then say \emph{bismôy} again
\end{itemize}

\begin{tcolorbox}[breakable,colback=teal!4,colframe=teal!35!black,title={Before you move on}]
What does \bn{বিস্ময়} mean? (\textquotedblleft{}Wonder; amazement\textquotedblright{}.) Say it once more.
\end{tcolorbox}
\section[\texorpdfstring{\bn{কৌতূহল} (koutūhôl)}{কৌতূহল (koutūhôl)}]{\bn{কৌতূহল} (koutūhôl) — curiosity}

\label{lesson:BN-C440-koutuhol}



\begin{quote}
Before the new one: say the Bengali for to get carried away, then the Bengali for wonder.
\end{quote}

\subsection*{You'll want to know: \bn{কৌতূহল}}
\textbf{\bn{কৌতূহল}} — \emph{koutūhôl} — \textquotedblleft{}curiosity\textquotedblright{}.

\begin{quote}
\bn{নতুন} \bn{বাক্সটা} \bn{দেখে} \bn{ছেলেটার} \bn{কৌতূহল} \bn{হল।}
\end{quote}

\begin{quote}
\emph{nôtun bāksôṭā dekhe chheleṭār koutūhôl hôlô.}
\end{quote}

\begin{quote}
Seeing the new box, the boy grew curious.
\end{quote}

\bn{কৌতূহলী}, curious.

\begin{cousinweb}
\textbf{\bn{কৌতূহল}} (\emph{koutūhôl}) is a tatsama, Sanskrit \emph{kautūhala}, \textquotedblleft{}curiosity, eager interest\textquotedblright{} (Monier-Williams); it shows two of the new vowel signs, \bn{ৌ} and \bn{ূ}, in one word.
\end{cousinweb}

\begin{grammarlens}[title={What you've built}]
One more word for describing what you lived and felt.
\end{grammarlens}

\subsection*{Your turn}
\begin{itemize}
\raggedright
  \item \emph{Say it:} \emph{koutūhôl}
  \item \emph{Say it:} \emph{koutūhôl}, once more
  \item \emph{Say it:} \emph{nôtun bāksôṭā dekhe chheleṭār koutūhôl hôlô.} — Seeing the new box, the boy grew curious.
  \item \emph{Recall:} say the Bengali for to get carried away, then the Bengali for wonder, then say \emph{koutūhôl} again
\end{itemize}

\begin{tcolorbox}[breakable,colback=teal!4,colframe=teal!35!black,title={Before you move on}]
What does \bn{কৌতূহল} mean? (\textquotedblleft{}Curiosity\textquotedblright{}.) Say it once more.
\end{tcolorbox}
\section[\texorpdfstring{\bn{মুগ্ধ} (mugdhô)}{মুগ্ধ (mugdhô)}]{\bn{মুগ্ধ} (mugdhô) — enchanted; spellbound}

\label{lesson:BN-C440-mugdho}



\begin{quote}
Before the new one: say the Bengali for wonder, then the Bengali for curiosity.
\end{quote}

\subsection*{You'll want to know: \bn{মুগ্ধ}}
\textbf{\bn{মুগ্ধ}} — \emph{mugdhô} — \textquotedblleft{}enchanted; spellbound\textquotedblright{}.

\begin{quote}
\bn{মীরার} \bn{গান} \bn{শুনে} \bn{আমি} \bn{মুগ্ধ} \bn{হলাম।}
\end{quote}

\begin{quote}
\emph{mīrār gān shune āmi mugdhô hôlām.}
\end{quote}

\begin{quote}
Hearing Mira sing, I was spellbound.
\end{quote}

\begin{cousinweb}
\textbf{\bn{মুগ্ধ}} (\emph{mugdhô}) is a tatsama, Sanskrit \emph{mugdha}, \textquotedblleft{}bewildered\textquotedblright{} and also \textquotedblleft{}lovely, charming\textquotedblright{} (Monier-Williams); Bengali keeps the sense of being held by something beautiful.
\end{cousinweb}

\begin{grammarlens}[title={What you've built}]
One more word for describing what you lived and felt.
\end{grammarlens}

\subsection*{Your turn}
\begin{itemize}
\raggedright
  \item \emph{Say it:} \emph{mugdhô}
  \item \emph{Say it:} \emph{mugdhô}, once more
  \item \emph{Say it:} \emph{mīrār gān shune āmi mugdhô hôlām.} — Hearing Mira sing, I was spellbound.
  \item \emph{Recall:} say the Bengali for wonder, then the Bengali for curiosity, then say \emph{mugdhô} again
\end{itemize}

\begin{tcolorbox}[breakable,colback=teal!4,colframe=teal!35!black,title={Before you move on}]
What does \bn{মুগ্ধ} mean? (\textquotedblleft{}Enchanted; spellbound\textquotedblright{}.) Say it once more.
\end{tcolorbox}
\section[\texorpdfstring{\bn{অবাক} \bn{হওয়া} (ôbāk hôoyā)}{অবাক হওয়া (ôbāk hôoyā)}]{\bn{অবাক} \bn{হওয়া} (ôbāk hôoyā) — to be surprised; to be amazed}

\label{lesson:BN-C440-obak-hooya}



\begin{quote}
Before the new one: say the Bengali for curiosity, then the Bengali for enchanted.
\end{quote}

\subsection*{You'll want to know: \bn{অবাক} \bn{হওয়া}}
\textbf{\bn{অবাক} \bn{হওয়া}} — \emph{ôbāk hôoyā} — \textquotedblleft{}to be surprised; to be amazed\textquotedblright{}.

\begin{quote}
\bn{বাবাকে} \bn{হঠাৎ} \bn{দেখে} \bn{আমি} \bn{অবাক} \bn{হলাম।}
\end{quote}

\begin{quote}
\emph{bābāke hôṭhāt dekhe āmi ôbāk hôlām.}
\end{quote}

\begin{quote}
I was surprised to see my father all of a sudden.
\end{quote}

Past (he/she …): \bn{তিনি} \bn{খবরটা} \bn{শুনে} \bn{অবাক} \bn{হলেন}, \emph{tini khôbôrṭā shune ôbāk hôlen}.

\begin{cousinweb}
\textbf{\bn{অবাক} \bn{হওয়া}} (\emph{ôbāk hôoyā}) is \emph{ôbāk}, \textquotedblleft{}speechless, dumbfounded\textquotedblright{}, which CDIAL lists for Bengali (CDIAL 882a), with \emph{hôoyā}, \textquotedblleft{}to be, to become\textquotedblright{} (CDIAL 9416): to be struck dumb.
\end{cousinweb}

\begin{grammarlens}[title={What you've built}]
One more word for describing what you lived and felt.
\end{grammarlens}

\subsection*{Your turn}
\begin{itemize}
\raggedright
  \item \emph{Say it:} \emph{ôbāk hôoyā}
  \item \emph{Say it:} \emph{ôbāk hôoyā}, once more
  \item \emph{Say it:} \emph{bābāke hôṭhāt dekhe āmi ôbāk hôlām.} — I was surprised to see my father all of a sudden.
  \item \emph{Recall:} say the Bengali for curiosity, then the Bengali for enchanted, then say \emph{ôbāk hôoyā} again
\end{itemize}

\begin{tcolorbox}[breakable,colback=teal!4,colframe=teal!35!black,title={Before you move on}]
What does \bn{অবাক} \bn{হওয়া} mean? (\textquotedblleft{}To be surprised; to be amazed\textquotedblright{}.) Say it once more.
\end{tcolorbox}
\section[\texorpdfstring{\bn{চমকানো} (chômkāno)}{চমকানো (chômkāno)}]{\bn{চমকানো} (chômkāno) — to be startled; to give a start}

\label{lesson:BN-C440-chomkano}



\begin{quote}
Before the new one: say the Bengali for enchanted, then the Bengali for to be surprised.
\end{quote}

\subsection*{You'll want to know: \bn{চমকানো}}
\textbf{\bn{চমকানো}} — \emph{chômkāno} — \textquotedblleft{}to be startled; to give a start\textquotedblright{}.

\begin{quote}
\bn{দরজায়} \bn{আওয়াজ} \bn{শুনে} \bn{মীরা} \bn{চমকে} \bn{গেল।}
\end{quote}

\begin{quote}
\emph{dôrjāy āoyāj shune mīrā chômke gelô.}
\end{quote}

\begin{quote}
Mira gave a start at a noise at the door.
\end{quote}

Past (he/she …): \bn{তিনি} \bn{চমকালেন}, \emph{tini chômkālen}.

Most often in its -\bn{এ} form: \bn{চমকে} \bn{যাওয়া} or \bn{চমকে} \bn{ওঠা}, to give a start.

\begin{cousinweb}
\textbf{\bn{চমকানো}} (\emph{chômkāno}) is listed by CDIAL for Bengali, as both glittering and being startled, under a reconstructed \emph{camakka}, \textquotedblleft{}a sudden movement\textquotedblright{} (CDIAL 4676).
\end{cousinweb}

\begin{grammarlens}[title={What you've built}]
One more word for describing what you lived and felt.
\end{grammarlens}

\subsection*{Your turn}
\begin{itemize}
\raggedright
  \item \emph{Say it:} \emph{chômkāno}
  \item \emph{Say it:} \emph{chômkāno}, once more
  \item \emph{Say it:} \emph{dôrjāy āoyāj shune mīrā chômke gelô.} — Mira gave a start at a noise at the door.
  \item \emph{Recall:} say the Bengali for enchanted, then the Bengali for to be surprised, then say \emph{chômkāno} again
\end{itemize}

\begin{tcolorbox}[breakable,colback=teal!4,colframe=teal!35!black,title={Before you move on}]
What does \bn{চমকানো} mean? (\textquotedblleft{}To be startled; to give a start\textquotedblright{}.) Say it once more.
\end{tcolorbox}
